% Préambule V2 : STYLE CLASSIQUE DE MONOGRAPHIE. Appelé par main_v2.tex.
%
% CE FICHIER NE CHANGE QUE L'APPARENCE. Les chapitres sont EXACTEMENT les mêmes
% fichiers que ceux de main.tex : aucun mot, aucun nombre, aucune structure ne
% diffère entre les deux versions. C'est la condition pour que la comparaison
% porte sur le style seul, et c'est aussi ce qui évite de maintenir deux textes.
%
% CE QUI CHANGE, ET POURQUOI
%   1. POLICE. Palatino (mathpazo) au lieu de Latin Modern. C'est une romaine de
%      la Renaissance revisitée : c'est elle qui donne l'aspect << texte ancien >>.
%      Sa hauteur d'x est grande, donc l'interligne est ouvert à 1,05, sans quoi
%      le gris de la page devient trop dense.
%   2. AUCUNE COULEUR. Tous les noms de couleurs du projet sont redéfinis en NOIR,
%      ce qui neutralise les \textcolor{navy} des chapitres sans les éditer. Les
%      liens restent cliquables mais s'impriment en noir.
%   3. LES ENCADRÉS DEVIENNENT DES FILETS. Les 66 blocs << clé >> et les 48 blocs
%      << attention >> perdent leur fond bleu et leur cadre arrondi. Ils gardent
%      leur distinction, mais par la TYPOGRAPHIE et non par la couleur : filets
%      haut et bas pour la synthèse, filet vertical épais pour l'avertissement.
%      C'est le dispositif des monographies imprimées, et il survit à la
%      photocopie comme à l'impression en noir et blanc.
%   4. TITRES. Petites capitales et filets, jamais de gras coloré. Les têtes de
%      paragraphe passent en italique : elles sont ici des phrases entières, et
%      une phrase entière en gras alourdit la page.
%   5. TITRES COURANTS. Chapitre en petites capitales en tête de page. Signal
%      classique, et utile sur un document de cette longueur.
%
% CE QUI NE CHANGE PAS, ET C'EST DÉLIBÉRÉ
%   LA GÉOMÉTRIE. Les marges restent à 2,4 cm. Une monographie classique aurait
%   une justification plus étroite, mais la lisibilité des cinquante figures est
%   calibrée sur CETTE largeur : \figover déborde à 1,14 fois la justification et
%   \figcle remplit la page. Rétrécir le bloc de texte rendrait les étiquettes
%   illisibles, ce qui est le défaut que la note de \figover documente. Le style
%   ne se paie pas en lisibilité des résultats.
%
%   LES CHIFFRES RESTENT ALIGNÉS (pas de chiffres elzéviriens). Palatino en
%   propose, et ils sont plus élégants dans un texte courant. Ils sont refusés
%   ici : ce document se lit pour ses nombres, et des chiffres de hauteur
%   inégale se comparent mal dans une colonne de tableau.

\usepackage[utf8]{inputenc}
\usepackage[T1]{fontenc}
\usepackage[sc]{mathpazo}          % Palatino, avec vraies petites capitales
\linespread{1.05}                  % obligatoire avec Palatino : grande hauteur d'x
\usepackage[french]{babel}
\usepackage[margin=2.4cm]{geometry}
\usepackage{microtype}

% Palatino est PLUS LARGE que Latin Modern a taille nominale egale, si bien que
% des paragraphes qui tenaient debordent. On autorise donc TeX a etirer les
% blancs entre mots en dernier recours, ce qui resorbe les debordements de
% PARAGRAPHE sans toucher une ligne des chapitres. Les debordements de TABLEAU,
% eux, ne s'y reduisent pas : une colonne trop large reste trop large, et cela
% se corrige dans le tableau ou pas du tout.
\setlength{\emergencystretch}{2.5em}
\usepackage{booktabs}
\usepackage{array}
\usepackage[table]{xcolor}
\usepackage{amsmath}
\usepackage{amssymb}
\usepackage{mathtools}
\usepackage{enumitem}
\usepackage[most]{tcolorbox}
\usepackage{tabularx}
\usepackage{longtable}
\usepackage{graphicx}
\usepackage{eurosym}
\usepackage{caption}
\usepackage{titlesec}
\usepackage{fancyhdr}
\usepackage{amsthm}
\usepackage[round,authoryear]{natbib}
\usepackage{hyperref}

% ---------------------------------------------------------------------------
% FIGURES : dispositif IDENTIQUE à la version 1, et il ne doit pas bouger.
% Voir la note du préambule de la version 1 pour le motif complet.
% ---------------------------------------------------------------------------
\newcommand{\figover}[1]{%
  \makebox[\textwidth][c]{\includegraphics[width=1.14\textwidth]{#1}}}

\newcommand{\figcle}[3]{%
  \begin{figure}[p]
    \centering
    \makebox[\textwidth][c]{%
      \includegraphics[width=1.10\textwidth,height=0.88\textheight,keepaspectratio]{#1}}
    \caption{#2}
    \label{#3}
  \end{figure}}

\graphicspath{{../vasicek_lab/figures/}{../cascade_qualitative/figures/}{../../outputs/}}

% ---------------------------------------------------------------------------
% ÉNONCÉS. Style classique : corps en italique pour les théorèmes, en romain
% pour les définitions, ce qui les distingue sans numérotation séparée.
% ---------------------------------------------------------------------------
\theoremstyle{plain}
\newtheorem{theorem}{Théorème}[chapter]
\newtheorem{proposition}[theorem]{Proposition}
\newtheorem{lemma}[theorem]{Lemme}
\theoremstyle{definition}
\newtheorem{definition}[theorem]{Définition}

\newcommand{\R}{\mathbb{R}}
\newcommand{\E}{\mathbb{E}}
\newcommand{\Prob}{\mathbb{P}}
\newcommand{\Var}{\mathrm{Var}}
% VaR et TVaR sont RESERVEES a la charge annuelle agregee L, c'est-a-dire a la mesure de
% capital. Le quantile de la severite d'un SINISTRE UNIQUE est un autre objet : il porte le
% symbole q, et sa moyenne de queue le symbole q barre. Ne pas les confondre : un niveau de
% 99,5 % n'a de sens reglementaire que sur un horizon ANNUEL, et une perte unique n'en porte
% aucun. Voir la table des notations et le chapitre socle.
\newcommand{\VaR}{\mathrm{VaR}}
\newcommand{\TVaR}{\mathrm{TVaR}}
\newcommand{\qsev}{\mathrm{q}}
\newcommand{\qbarsev}{\overline{\mathrm{q}}}
\newcommand{\dd}{\mathrm{d}}
\newcommand{\indic}{\mathbf{1}}
\providecommand{\up}[1]{\textsuperscript{#1}}

% ---------------------------------------------------------------------------
% PALETTE : tous les noms de la version 1 existent encore, et ils valent NOIR.
% C'est ce qui rend les chapitres réutilisables sans une seule édition : les
% cinq \textcolor{navy} de la table des notations s'impriment en noir, et les
% cellules colorées des tables de synthèse deviennent blanches.
% ---------------------------------------------------------------------------
\definecolor{navy}{gray}{0}
\definecolor{bluemid}{gray}{0}
\definecolor{lightblue}{gray}{1}
\definecolor{accent}{gray}{0}
\definecolor{green}{gray}{0}
\definecolor{lightgreen}{gray}{1}
\definecolor{softgreen}{gray}{1}
\definecolor{softamber}{gray}{1}

\hypersetup{colorlinks=true, linkcolor=black, urlcolor=black, citecolor=black,
  pdftitle={Quantification du SCR lie a la non-conformite DORA par la cascade dirigee},
  pdfauthor={Kelian Kaddouri}}

% ---------------------------------------------------------------------------
% LES DEUX ENCADRÉS, EN TYPOGRAPHIE PLUTÔT QU'EN COULEUR.
% Ils restent des environnements tcolorbox pour une raison pratique : ils
% doivent pouvoir se COUPER entre deux pages, ce qu'un simple \hrule autour
% d'un bloc ne sait pas faire proprement sur des blocs de dix lignes.
%   cle       : filets haut et bas, pleine justification. C'est le << résumé
%               encadré >> des monographies : il arrête l'oeil sans le colorer.
%   attention : filet vertical épais et retrait. Distinction lisible en noir et
%               blanc, et hiérarchie claire, l'avertissement étant en retrait
%               du fil du texte quand la synthèse l'interrompt.
% ---------------------------------------------------------------------------
% NOTE DE MISE EN OEUVRE, ET ELLE A COUTE UN ESSAI. Mettre boxrule a 0pt puis
% ne relever que toprule et bottomrule NE SUFFIT PAS : le squelette << enhanced >>
% trace quand meme un filet de cheveu sur les quatre cotes, et les deux encadres
% se retrouvent tous deux cadres, donc indistinguables. L'idiome correct est
% << frame hidden >> pour supprimer le cadre, puis des << borderline >> pour ne
% dessiner que les filets voulus. Verifie sur le PDF rendu, pas sur la source.
\newtcolorbox{cle}[1][]{%
  enhanced, breakable, sharp corners,
  colback=white, frame hidden,
  borderline north={0.9pt}{0pt}{black},
  borderline south={0.9pt}{0pt}{black},
  left=0pt, right=0pt, top=8pt, bottom=8pt,
  before skip=12pt, after skip=12pt, #1}

\newtcolorbox{attention}[1][]{%
  enhanced, breakable, sharp corners,
  colback=white, frame hidden,
  borderline west={1.8pt}{0pt}{black},
  left=12pt, right=0pt, top=4pt, bottom=4pt,
  before skip=12pt, after skip=12pt, #1}

\newtcolorbox{migration}[1][]{%
  enhanced, breakable, sharp corners, colback=white, colframe=black,
  boxrule=0.4pt, left=8pt, right=8pt, top=5pt, bottom=5pt,
  before skip=8pt, after skip=8pt, #1}
\newcommand{\areprendre}[1]{\begin{migration}\footnotesize\textit{[À migrer]}~#1\end{migration}}

% ---------------------------------------------------------------------------
% VERDICTS DES TABLES DE SYNTHÈSE. La version 1 les distingue par la couleur du
% fond de cellule. Ici c'est le mot qui porte, en petites capitales : les trois
% verdicts se lisent aussi bien, et la table survit à une impression en noir.
% \cascell devient neutre : la ligne qu'elle marquait porte déjà son gras.
% ---------------------------------------------------------------------------
\newcommand{\rob}{\textsc{robuste}}
\newcommand{\feat}{\textsc{voulu}}
\newcommand{\ass}{\textsc{assumé}}
\newcommand{\oui}{$\checkmark$}
\newcommand{\non}{$\times$}
\newcommand{\cascell}{}

% ---------------------------------------------------------------------------
% TITRES. Petites capitales et filets. Aucun gras coloré.
% L'espacement resserré de la version 1 est CONSERVÉ : il économise environ un
% quart de page par chapitre, et le style n'a pas de raison de le rendre.
% ---------------------------------------------------------------------------
\titleformat{\part}[display]
  {\normalfont\filcenter}
  {\large\scshape\partname~\thepart}
  {16pt}
  {\Huge\scshape}
\titleformat{\chapter}[display]
  {\normalfont}
  {\normalsize\scshape\chaptertitlename~\thechapter}
  {8pt}
  {\huge}
  [{\vspace{6pt}\titlerule[0.8pt]}]
\titlespacing*{\chapter}{0pt}{-10pt}{22pt}
\titleformat{\section}{\normalfont\large\bfseries}{\thesection}{0.6em}{}
\titleformat{\subsection}{\normalfont\normalsize\itshape}{\thesubsection}{0.5em}{}
\titleformat{\paragraph}[runin]{\normalfont\normalsize\itshape}{}{0pt}{}
\titlespacing*{\paragraph}{0pt}{9pt}{0.55em}

\captionsetup{font=small,labelfont=sc,labelsep=period,skip=6pt}

% Titres courants : chapitre en petites capitales, filet fin, folio à droite.
\pagestyle{fancy}
\fancyhf{}
\fancyhead[L]{\small\scshape\nouppercase{\leftmark}}
\fancyhead[R]{\small\thepage}
\renewcommand{\headrulewidth}{0.4pt}
\fancypagestyle{plain}{\fancyhf{}\fancyfoot[C]{\small\thepage}%
  \renewcommand{\headrulewidth}{0pt}}

% ---------------------------------------------------------------------------
% PAGE DE TITRE classique : filets, petites capitales, pas de couleur.
% ---------------------------------------------------------------------------
\title{%
  \vspace{-0.8cm}
  \rule{\textwidth}{1pt}\\[10pt]
  {\Large\scshape Quantification du SCR lié à la non-conformité\\[4pt]
   au règlement DORA}\\[10pt]
  \rule{\textwidth}{0.4pt}\\[12pt]
  {\large\itshape Une cascade dirigée entre les cinq piliers,\\
   dont l'état de conformité pilote le capital}}
\author{\normalsize\scshape Kélian Kaddouri \\[6pt]
  \normalfont\small ENSAE / Nexialog Consulting \\
  \small Tuteurs : Hugo Rapior, Caroline Hillairet}
\date{\small\today}
